% from ega3-1-fr.tex
\oldpage[III]{127}

(où les flèches horizontales sont la multiplication par $x$, et les
flèches verticales les injections canoniques) est commutatif~; donc le
diagramme correspondant
\[
\xymatrix{
H^n(X,\mathfrak J^{k+1}\mathcal F)\ar[r]^{\mu_{x,m}}\ar[d]&
H^n(X,\mathfrak J^{k+m+1}\mathcal F)\ar[d]\\
H^n(X,\mathcal F)\ar[r]_{\mu_{x,0}}&H^n(X,\mathcal F)
}
\]
est commutatif, ce qui, compte tenu de l'interprétation de $\mathrm R_k$
comme image de
$H^n(X,\mathfrak J^{k+1}\mathcal F)\to H^n(X,\mathcal F)$, prouve
(\hyperref[III.4.1.7.10-fr]{4.1.7.10}) et montre en outre que le $S$-module
gradué $\mathrm R=\bigoplus_{k\geq0}\mathrm R_k$ est un quotient du
sous-$S$-module
$\mathrm M=\bigoplus_{k\geq0}H^n(X,\mathfrak J^{k+1}\mathcal F)$ de
$\mathrm E$~; la remarque faite plus haut montre donc que $\mathrm R$ est
un $S$-module de type fini, ce qui équivaut à l'assertion de
(\hyperref[III.4.1.7.9-fr]{4.1.7.9}) (Bourbaki, \emph{Alg. comm.},
chap. III, §~3, n\textsuperscript{o}~1, th.~1).
\end{proof}

\begin{env}[4.1.7.12]
\label{III.4.1.7.12-fr}
Considérons maintenant le $S$-module gradué
\[
\mathrm N=\bigoplus_{k\geq0}H^{n+1}(X,\mathfrak J^{k+1}\mathcal F)
\]
défini comme dans (4.7.1.8)~; c'est encore un $S$-module de type fini en
vertu de (\hyperref[III.3.3.2-fr]{3.3.2})~; on a
$\mathrm Q_k\subset\mathrm N_k$ pour tout $k$ par
(\hyperref[III.4.1.7.5-fr]{4.1.7.5}) et le diagramme
(\hyperref[III.4.1.7.11-fr]{4.1.7.11}) où on remplace $n$ par $n+1$,
montre que
$S_m\mathrm Q_k=\mathfrak J^m\mathrm Q_k\subset\mathrm Q_{k+m}$.
Autrement dit, $\mathrm Q$ est un sous-$S$-module gradué de $\mathrm N$,
et par suite est de type fini.
\end{env}

\begin{env}[4.1.7.13]
\label{III.4.1.7.13-fr}
Désignons par $\alpha_m$ l'injection canonique
$\mathfrak J^m\to A$, qu'on peut écrire $S_m\to S_0$. Comme
$\mathfrak J^{k+1}\mathcal F_k=0$, le $A$-module
$H^n(X,\mathcal F_k)$ est annulé par $\mathfrak J^{k+1}$~; comme
$\mathrm Q_k$ est l'image du $A$-homomorphisme
$H^n(X,\mathcal F_k)\to H^{n+1}(X,\mathfrak J^{k+1}\mathcal F)$,
$\mathrm Q_k$, en tant que $A$-module, est aussi annulé par
$\mathfrak J^{k+1}$~; cela signifie encore que, dans le $S$-module
$\mathrm Q$, on a
\[
\alpha_{k+1}(S_{k+1})\mathrm Q_k=0.
\tag{4.1.7.14}\label{III.4.1.7.14-fr}
\]
Comme $\mathrm Q$ est un $S$-module de type fini, il existe un entier
$k_0$ et un entier $h$ tels que
$\mathrm Q_{k+h}=S_h\mathrm Q_k$ pour $k\geq k_0$ (II, 2.1.6, (ii))~;
on déduit de cette relation et de
(\hyperref[III.4.1.7.14-fr]{4.1.7.14}) qu'il existe un entier $r>0$ tel
que
\[
\alpha_r(S_r)\mathrm Q=0.
\tag{4.1.7.15}\label{III.4.1.7.15-fr}
\]
\end{env}

\begin{env}[4.1.7.16]
\label{III.4.1.7.16-fr}
Notons maintenant que l'injection canonique
$\mathfrak J^{k+m}\mathcal F\to\mathfrak J^k\mathcal F$ donne en passant
à la cohomologie un $A$-homomorphisme
\[
\nu_m:H^{n+1}(X,\mathfrak J^{k+m}\mathcal F)
\longrightarrow H^{n+1}(X,\mathfrak J^k\mathcal F)
\tag{4.1.7.17}\label{III.4.1.7.17-fr}
\]
et, pour tout $x\in\mathfrak J^m$, on a évidemment la factorisation
\[
\xymatrix@C=4.5em{
H^{n+1}(X,\mathfrak J^k\mathcal F)\ar[r]^{\mu_{x,m}}&
H^{n+1}(X,\mathfrak J^{k+m}\mathcal F)\ar[r]^{\nu_m}&
H^{n+1}(X,\mathfrak J^k\mathcal F)
}
\tag{4.1.7.18}\label{III.4.1.7.18-fr}
\]
dont le composé est $\mu_{x,0}$.
\end{env}

